Comparamos sección por sección cinco algoritmos publicados para \textsc{Minimum Vertex Cover}, uno o dos por paradigma (voraz, programación dinámica, divide y vencerás y aleatorizado), y con esa comparación construimos \hyvc. El híbrido reduce el grafo con reglas exactas, lo divide en componentes, despacha a programación dinámica las componentes cuya tabla es barata y ramifica sobre el resto con poda por cotas. Siempre devuelve una cubierta válida y, si termina, óptima.

En 46 instancias de seis familias, \hyvc{} certificó el óptimo en 42, frente a 20 del divide y vencerás de Weigt y Hartmann y 17 de la programación dinámica sobre descomposiciones en árbol. Los dos algoritmos base solo funcionan en su régimen; \hyvc{} cubre la unión de ambos y además resuelve familias que ninguno cierra. La ablación muestra que la ganancia no se reparte por igual: casi todo el efecto viene de las reglas de reducción y, en grafos de ancho de árbol pequeño, de la programación dinámica. La división dinámica en componentes y la cota superior Monte Carlo, que en la comparación teórica parecían prometedoras, no tuvieron un efecto medible en estas instancias.

El desarrollo también dejó una lección de diseño: el criterio para despachar a programación dinámica debe ser un \emph{costo} estimado, y no un umbral de ancho. Con un umbral fijo, \hyvc{} no despachaba grafos que la PD resolvía en centésimas de segundo; con un único tope de costo, gastaba segundos en tablas que la ramificación resolvía en milisegundos.

\paragraph{Trabajo futuro.} (i) Añadir una fase de búsqueda local (por ejemplo, la de FastVC o NuMVC) para que la respuesta \emph{anytime} mejore cuando la búsqueda exacta no termina; (ii) cotas inferiores más fuertes, como la relajación LP o una cobertura por cliques~\citep{akiba2016branch}, para cerrar las instancias densas y de ley de potencias; (iii) reemplazar la descomposición por grado mínimo por una mejor (p.\,ej. la heurística de \emph{min-fill} o un solucionador de treewidth); (iv) ajustar $\gamma$ automáticamente, ya que valores $\gamma\ge 8$ resultaron claramente mejores que el $\gamma=2$ usado; y (v) explorar la ramificación aleatorizada de Clinch et al.~\citep{clinch2025faster} como alternativa Las Vegas a la ramificación determinista.
